% !TEX program = platex
% 数学ノート集・共通テンプレート v2（積分論2026年度第1部の版面を参考）
% A5判。章内で定義・定理・例・注意等の番号を共通化。
% コンパイル： platex -interaction=nonstopmode math_note_template_v2.tex (2回)
%              dvipdfmx math_note_template_v2.dvi
\documentclass[a5paper,10pt,onecolumn]{jsarticle}
\usepackage[dvipdfmx,top=22mm,bottom=20mm,inner=19mm,outer=18mm,headheight=13pt,headsep=6mm]{geometry}
\usepackage{amsmath,amssymb,amsthm,mathtools}
\usepackage{enumitem,fancyhdr,titlesec}
\usepackage[dvipdfmx,bookmarks=true,colorlinks=false]{hyperref}
\usepackage{pxjahyper}

% ----- 差し替える文書情報 -----
\newcommand{\NoteTitle}{D加群のde Rham複体と解函手}
\newcommand{\NoteSeries}{数学ノート集}

% ----- 文字・組版 -----
% 参考PDFの質素な明朝＋Computer Modernの印象を引き継ぎながら，
% A5判で窮屈にならないよう本文の行間を少し確保する．
\linespread{1.13}
\setlength{\parindent}{1zw}
\setlength{\parskip}{0pt}
\setlength{\abovedisplayskip}{6pt plus 2pt minus 2pt}
\setlength{\belowdisplayskip}{6pt plus 2pt minus 2pt}
\setlength{\abovedisplayshortskip}{3pt plus 2pt minus 1pt}
\setlength{\belowdisplayshortskip}{4pt plus 2pt minus 1pt}
\setlist[enumerate]{topsep=3pt,itemsep=1pt,parsep=0pt,leftmargin=2.4zw}
\titleformat{\section}{\normalfont\normalsize\bfseries}{\thesection}{0.8zw}{}
\titlespacing*{\section}{0pt}{2.0ex plus .4ex}{.9ex}
\titleformat{\subsection}{\normalfont\normalsize\bfseries}{\thesubsection}{0.7zw}{}
\titlespacing*{\subsection}{0pt}{1.5ex plus .3ex}{.6ex}
\raggedbottom

% ----- 柱とページ番号：ページ上部、罫線なし -----
\fancypagestyle{mathnote}{%
  \fancyhf{}%
  \fancyhead[L]{\footnotesize \NoteSeries\ （\NoteTitle）}%
  \fancyhead[R]{\footnotesize \thepage}%
  \renewcommand{\headrulewidth}{0pt}%
  \renewcommand{\footrulewidth}{0pt}%
}
\pagestyle{mathnote}

% ----- 節内通し番号の定義・定理群 -----
\newtheoremstyle{mathnotetheorem}
 {4.5pt}{4.5pt}{\normalfont}{}%
 {\bfseries}{}{0.65zw}%
 {\thmname{#1}\thmnumber{ #2}\thmnote{（#3）}}
\theoremstyle{mathnotetheorem}
\newtheorem{theorem}{定理}[section]
\newtheorem{definition}[theorem]{定義}
\newtheorem{lemma}[theorem]{補題}
\newtheorem{proposition}[theorem]{命題}
\newtheorem{corollary}[theorem]{系}
\newtheorem{example}[theorem]{例}
\newtheorem{remark}[theorem]{注意}
\renewcommand{\proofname}{証明}
\renewcommand{\qedsymbol}{$\square$}
\makeatletter
\renewenvironment{proof}[1][\proofname]{%
  \par\pushQED{\qed}\normalfont\topsep4\p@\@plus2\p@\relax
  \trivlist\item[\hskip\labelsep\bfseries #1]\ignorespaces
}{\popQED\endtrivlist\@endpefalse}
\makeatother

% ----- 使うノートで適宜拡張する記号 -----
\newcommand{\R}{\mathbb{R}}
\newcommand{\norm}[1]{\lVert #1\rVert}
\newcommand{\abs}[1]{\lvert #1\rvert}

\begin{document}
\section*{読了後にできること}
本稿を読むための前提は，代数的de Rham複体，Weyl代数，導関数．
本稿の具体的な到達課題は，多項式のde Rham複体のコホモロジーを計算する．

\section{de Rham複体の代数的定義}
$k$ は標数零の体とし，$A_n$ は $x_i,\partial_i$ のWeyl代数とする．
$M$ を左 $A_n$-加群とする．
\begin{definition}
$\Omega_k^p=\bigwedge^p k^n$ の基底を $dx_i$ の外積で表し，
$\operatorname{DR}^p(M)=M\otimes_k\Omega_k^p$ とする．
微分を
\[
 d(m\otimes\omega)=\sum_{i=1}^n\partial_i m\otimes(dx_i\wedge\omega)
\]
で定める．
\end{definition}
\begin{proposition}
$d^2=0$ である．従って $(\operatorname{DR}^{\bullet}(M),d)$ は複体である．
\end{proposition}
\begin{proof}
$\partial_i\partial_j=\partial_j\partial_i$ に対し
$dx_i\wedge dx_j=-dx_j\wedge dx_i$ である．
$d^2$ の二重和は添字 $i,j$ を入れ替えると符号が変わり，対ごとに消える．
\end{proof}
\section{多項式環のPoincar\'e補題}
\begin{theorem}
$M=k[x_1,\ldots,x_n]$ を標準 $A_n$-加群とする．
このとき，そのde Rhamコホモロジーは
\[
 H^p(\operatorname{DR}(M))=
 \begin{cases}k&(p=0),\\0&(p>0)\end{cases}
\]
である．
\end{theorem}
\begin{proof}
多項式係数の微分形式に総次数
$N=$(係数の多項式次数)+(微分形式次数) を与える．
Eulerベクトル場 $E=\sum_i x_i\partial_i$ による内部積 $\iota_E$ を用いる．
同次数 $N$ の斉次形式に対し，直接各単項式で確かめると
\[
 d\iota_E+\iota_Ed=N\operatorname{id}
\]
となる．$N>0$ で $d\omega=0$ なら
$\omega=d(N^{-1}\iota_E\omega)$ である．
標数零なので $N$ で割れる．$N=0$ の形式は定数の零次形式のみであり，
これは正確に $H^0=k$ に寄与する．
\end{proof}
\section{解函手と具体例}
\begin{definition}
微分方程式 $P\cdot u=0$ の解を $A_1$-線形写像
$A_1/A_1P\to\mathcal O(U)$ として扱う．
ここで $U\subset\mathbb C$ は開集合，$\mathcal O(U)$ は正則関数で
通常の微分作用を持つものとする．
\end{definition}
\begin{proposition}
$U$ を連結開集合，$\lambda\in\mathbb C$ とする．
$P=\partial-\lambda$ に対する上記の解空間は一次元であり，
$e^{\lambda z}$ を基底とする．
\end{proposition}
\begin{proof}
準同型は $1+A_1P$ の像 $u\in\mathcal O(U)$ で決まり，
$u'=\lambda u$ を満たす．積の微分により
$(e^{-\lambda z}u)'=0$ であり，連結性からこれは定数である．
逆に任意の定数倍は条件を満たす．
\end{proof}
\begin{remark}
解層の導来函手とRiemann--Hilbert対応は本稿では構成していない．
ここまでの複体および具体例は，そのための初等的準備である．
\end{remark}


\section{正則特異点の一次元モデルとモノドロミー}
\begin{proposition}[一次方程式の局所解とモノドロミー]
$\lambda\in\mathbb C$ とする．$U=\mathbb C^*$ 上の方程式
\[
 z\frac{du}{dz}=\lambda u
\]
は各単連結開部分集合 $V\subset U$ 上で一次元の正則解空間を持つ．
解を原点の周りに一周解析接続すると $e^{2\pi i\lambda}$ 倍になる．
\end{proposition}
\begin{proof}
単連結 $V\subset\mathbb C^*$ には正則な対数 $\ell(z)$ が存在する．
$u(z)=c\exp(\lambda\ell(z))$ は表示した方程式を満たす．
逆に任意の解 $u$ について
$\frac{d}{dz}(e^{-\lambda\ell(z)}u(z))=0$ であるから定数倍以外の解はない．
対数を原点の周りに一周解析接続すると $\ell$ が $2\pi i$ 増加するため，
解は指定された因子だけ変化する．
\end{proof}
\begin{theorem}[一次元局所系の分類]
$\mathbb C^*$ 上の複素数一次元局所系の同型類は
$\mathbb C^*$ の元 $c\ne0$ と一対一に対応する．
前命題の方程式の解局所系は $c=e^{2\pi i\lambda}$ に対応する．
\end{theorem}
\begin{proof}
$\mathbb C^*$ の普遍被覆を $p:\mathbb C\to\mathbb C^*$，
$p(w)=e^w$ とする．局所系を引き戻すと，単連結な $\mathbb C$ 上では
一つの一次元ベクトル空間 $V$ に定値な局所系になる．
被覆変換 $w\mapsto w+2\pi i$ は $V$ 上の可逆線形写像を与える．
一次元なのでこれは非零スカラー $c$ である．
逆に，任意の $c\in\mathbb C^*$ に対し
$\mathbb C\times\mathbb C$ を
$(w,v)\sim(w+2\pi i,c^{-1}v)$ で割ると局所系が得られる．
二つの対応は互いに逆であり，前命題から最後の記述が従う．
\end{proof}
\begin{remark}
この分類は $\mathbb C^*$ 上の階数1の場合に限ったものであり，
一般のRiemann--Hilbert対応の証明を代用するものではない．
\end{remark}
\section{一定係数の正則特異接続}
\begin{theorem}
$A\in M_r(\mathbb C)$ を固定する．$\mathbb C^*$ 上の連立方程式
$z\,u'(z)=Au(z)$ の局所正則解空間は常に $r$ 次元である．
原点を一周するモノドロミーは，基底の選択を除いて
$e^{2\pi iA}$ に等しい．
\end{theorem}
\begin{proof}
単連結な開集合 $V\subset\mathbb C^*$ 上で対数 $\ell$ を選び，
$Y(z)=\exp(A\ell(z))$ と置く．
行列 $A$ と $\exp(A\ell)$ は可換なので
$Y'(z)=z^{-1}AY(z)$ が成り立つ．
$\det Y(z)=\exp(\operatorname{tr}A\,\ell(z))\ne0$ だから，
$Y$ は基本解行列である．任意の解について
$(Y^{-1}u)'=0$ となるので，局所解は $u(z)=Y(z)c$（$c\in\mathbb C^r$）に限られる．
対数を原点の周りに一周解析接続すると $\ell$ は $2\pi i$ 増すので
$Y$ は右から $e^{2\pi iA}$ を掛けたものに移る．
従って基底を $Y$ の列とすればモノドロミーはその行列で表される．
\end{proof}

\section{非特異な平坦接続と局所系の同値}
この節では $X$ を複素多様体とし，有限階数の正則ベクトル束
$E\to X$ に正則接続 $\nabla$ が与えられた状況を考える．
平坦とは $\nabla^2=0$，すなわち曲率が零であることをいう．
\begin{lemma}[正則なFrobenius型可積分系]\label{lem:flatmatrix}
原点を含む小さい多重円板 $U\subset\mathbb C^n$ で，
$A_1,\ldots,A_n:U\to M_r(\mathbb C)$ を正則とする．
\[
 \partial_iA_j-\partial_jA_i=[A_i,A_j]
 \qquad(1\le i,j\le n)
\]
を満たすなら，原点の近傍に正則行列関数
$Y$ が一意に存在して
\[
 \partial_iY=A_iY,\qquad Y(0)=I_r
\]
を満たす．しかも $\det Y$ は零にならない．
\end{lemma}
\begin{proof}
座標を一つずつ増やして構成する．
まず直線 $z_2=\cdots=z_n=0$ 上で
$\partial_1Y=A_1Y$，$Y(0)=I_r$ を解く．
存在一意性と初期値・係数に関する正則性は，線形常微分方程式の
Picard逐次近似を小さい円板上で行えば従う．
次に $z_1,z_2$ に関する部分空間へ延長する．
既に決めた $Y(z_1,0)$ を初期値として
$\partial_2Y=A_2Y$ を解けば，$Y$ は二変数について正則になる．
$E_1=\partial_1Y-A_1Y$ と置くと，交換関係から
\[
 \begin{aligned}
 \partial_2E_1
 &=\partial_1(A_2Y)-(\partial_2A_1)Y-A_1A_2Y\\
 &=A_2E_1+
 (\partial_1A_2-\partial_2A_1+A_2A_1-A_1A_2)Y
 =A_2E_1.
 \end{aligned}
\]
$z_2=0$ 上では $E_1=0$ だから，線形常微分方程式の
一意性により $E_1$ は恒等的に零となる．
同様に $z_3,\ldots,z_n$ の方向へ順に延長する．
既に満たしていた $\partial_jY=A_jY$ は，新しい方向の
初期値問題の一意性と上記の計算によって保存される．
これで存在が得られ，同じ一意性の議論を各座標方向に適用すれば
$Y(0)=I_r$ からの一意性も従う．
$Y$ の各列は微分系の解である．行列式については
$\partial_i\det Y=(\operatorname{tr}A_i)\det Y$ だから，
原点で $\det Y=1$ であることと各方向の一意性から
$\det Y\ne0$ が近傍で成立する．
\end{proof}
\begin{theorem}[非特異なRiemann--Hilbert対応]
複素多様体 $X$ 上の有限階数正則ベクトル束で平坦な正則接続を持つものの
圏と，有限階数の複素局所系の圏は同値である．
対応は，平坦接続からその水平切断の局所系を作る操作と，
局所系 $\mathcal L$ から
$E=\mathcal O_X\otimes_{\mathbb C}\mathcal L$ を作って
微分を接続とする操作で与えられる．
\end{theorem}
\begin{proof}
平坦接続 $(E,\nabla)$ を局所自明化し，
$\nabla_i=\partial_i-A_i$ と書く．
平坦性は前補題の交換関係に等しい．従って小さい多重円板上で
可逆な正則基本解行列 $Y$ が存在する．
その水平切断は $Y(z)c$（$c\in\mathbb C^r$）に限られるから，
水平切断の層 $\mathcal L=\ker\nabla$ は局所的に定数層である．
評価写像
\[
 \mathcal O_X\otimes_{\mathbb C}\mathcal L
 \longrightarrow E,\qquad f\otimes s\longmapsto fs
\]
は小さい多重円板で $Y$ を用いた自明化に一致するため，
正則ベクトル束と平坦接続の同型を与える．

逆に局所系 $\mathcal L$ の局所定値基底を取ると，
$\mathcal O_X\otimes_{\mathbb C}\mathcal L$ は
正則ベクトル束となり，$\nabla(f\otimes s)=df\otimes s$ と定められる．
異なる局所定値基底の遷移行列は局所定数だから，この接続の定義は
重なり上で一致する．$d^2=0$ より接続は平坦であり，
水平切断はちょうど $\mathcal L$ である．
束と接続を保つ射は水平切断を保ち，
局所系の射はテンソル積に拡張される．
両構成に伴う自然な同型を上に構成したので，圏同値が従う．
\end{proof}
\begin{remark}
本節では特異点のない有限階数接続に対する対応を証明した．
正則ホロノミック $D$-加群と構成可能複体の導来圏との一般の
Riemann--Hilbert対応には，特異点での延長と中間拡張，
perverse t-構造，双対性など別の論証が必要である．
\end{remark}

\section{任意の有限次元モノドロミーの対数接続による実現}
\begin{lemma}[可逆複素行列の対数]
任意の $M\in\operatorname{GL}_r(\mathbb C)$ に対し，
ある $B\in M_r(\mathbb C)$ が存在して
$e^{2\pi iB}=M$ を満たす．
\end{lemma}
\begin{proof}
$M$ をJordan標準形 $M=PJP^{-1}$ に移す．
各Jordanブロックは $J_\lambda=\lambda(I+N)$
（$\lambda\ne0$，$N$ は冪零）と書ける．
$e^{\ell_\lambda}=\lambda$ を満たす対数
$\ell_\lambda\in\mathbb C$ を選び，
\[
 L_\lambda=\ell_\lambda I
 +\sum_{j=1}^{r-1}\frac{(-1)^{j+1}}jN^j
\]
と置く．$N^r=0$ なので和は有限である．
形式的冪級数の恒等式
$\exp(\log(1+z))=1+z$ を冪零行列 $N$ に代入すると
$e^{L_\lambda}=\lambda(I+N)=J_\lambda$ を得る．
各ブロックの $L_\lambda$ を直和した行列を $L$ とすると
$B=(2\pi i)^{-1}PLP^{-1}$ は $e^{2\pi iB}=M$ を満たす．
\end{proof}
\begin{theorem}[対数型接続によるモノドロミーの実現]
穴あき円板 $\Delta^*=\{z\in\mathbb C:0<|z|<1\}$ とし，
$M\in\operatorname{GL}_r(\mathbb C)$ を任意とする．
このとき $\Delta^*$ 上の階数 $r$ の正則平坦接続で，
$z=0$ に高々一次極を持つ対数型接続が存在して，
その正向き一周のモノドロミーが $M$ に共役となる．
\end{theorem}
\begin{proof}
前補題により $e^{2\pi iB}=M$ を満たす行列 $B$ を取る．
自明束 $\Delta^*\times\mathbb C^r$ 上の接続を
\[
 \nabla=d-B\frac{dz}{z}
\]
と定義する．係数行列 $B\,dz/z$ は $z=0$ で一次極を持つ．
基底空間が複素一次元なので曲率
$d(B\,dz/z)-(B\,dz/z)\wedge(B\,dz/z)$ は零であり，
この接続は平坦である．
単連結な小領域で対数の枝 $\operatorname{Log}z$ を固定すると，
水平切断の基本行列は
$Y(z)=\exp(B\operatorname{Log}z)$ である．
実際 $dY=(B\,dz/z)Y$ が成立する．
原点を正向きに一周すると
$\operatorname{Log}z$ は $2\pi i$ 増えるから，
$Y$ は $Y\,e^{2\pi iB}=YM$ に変わる．
基点での基底の選び方によって共役が生じるので，
モノドロミーの共役類が $M$ のそれに一致する．
\end{proof}
\begin{remark}
任意の有限次元モノドロミー表現
$\pi_1(\Delta^*)\simeq\mathbb Z\to\operatorname{GL}_r(\mathbb C)$
を対数型接続で実現できることは示した．
正則ホロノミックD加群全体の導来Riemann--Hilbert対応および
perverse sheafの理論はこの主張よりはるかに一般的であり，
本節では証明していない．
\end{remark}


\section{到達課題の完全解答}
\noindent\textbf{一次元の代数的Poincaré補題．}
\begin{proposition}
標数零の体 $k$ 上のアフィン直線における
代数的de Rham複体
\[
0\longrightarrow k[x]
\xrightarrow{d}k[x]\,dx\longrightarrow0
\]
のコホモロジーは
\[
H^0_{\mathrm{dR}}(\mathbb A^1_k)=k,\qquad
H^1_{\mathrm{dR}}(\mathbb A^1_k)=0
\]
である．
\end{proposition}
\begin{proof}
$d(f)=f'(x)\,dx$ なので核は定数多項式である．
また任意の $g(x)=\sum_{j=0}^ma_jx^j$ に対し
\[
F(x)=\sum_{j=0}^m\frac{a_j}{j+1}x^{j+1}
\]
は $F'=g$ を満たす．
標数零なので各 $j+1$ は $k$ で可逆である．
よって $d$ は一次形式全体へ全射であり，
一次コホモロジーは零である．
\end{proof}
\noindent\textbf{到達確認．}
$g(x)=3x^2+2$ の原始関数と対応する
一次形式のコホモロジー類を求めよ．
\par\noindent\textbf{解答．}
$F(x)=x^3+2x$ で $dF=g\,dx$．
従って $g\,dx$ の一次コホモロジー類は零である．
\end{document}
